\documentclass[aps,prl,twocolumn,nofootinbib,floatfix]{revtex4-2}

\usepackage{amsmath,amssymb,amsfonts}
\usepackage{graphicx}
\usepackage{hyperref}
\usepackage{cleveref}

\title{The Source Code of the Universe: Formal Deduction from Discrete Information to Cosmic Density}

\author{Zhang Jun}
\affiliation{New Beginning Universe Research Group (Independent Researcher)}
\email{cnjun939@163.com}
\orcid{0009-0004-9803-3237}

\begin{document}

\maketitle

\begin{abstract}
We present the first fully formalized deductive chain from information-theoretic axioms to an observable cosmological parameter, with zero free parameters. A characteristic constant $\theta = 1/(2+2\cos(2\pi/7)) \approx 0.308$ emerges uniquely from the cyclic group $\mathbb{Z}_7$, matching Planck 2018 $\Omega_m = 0.311$ within $\sim 1\%$. Among all primes $p$, only $p=7$ yields $\theta(p)$ within the structure formation window $(0.28, 0.33)$---establishing that the observed matter density is not a fitted parameter but a forced constraint of algebraic structure. The complete proof is machine-verified in Lean 4 (2196 compilation tasks, 0 errors), and the framework reframes observers as internal nodes of the causal lattice: structure$=7$ is the necessary condition for observers like us to exist.
\end{abstract}

\section{Introduction}

The reconciliation of quantum mechanics and general relativity remains unresolved. Traditional approaches---string theory, loop quantum gravity, causal set theory---each face fundamental challenges: no unique vacuum, no complete semiclassical limit, or no direct observational anchor. CSQIT takes a fundamentally different approach: rather than quantizing geometry or postulating extra dimensions, we start from a minimal set of information-theoretic axioms and use machine-verifiable formalization to derive quantitative consequences. This represents a third path---axiomatic deduction driven by formal proof---which yields a complete deductive chain from axioms to an observable cosmological constant, with zero free parameters.

\section{Axiomatic System and Formalization}

The CSQIT framework is built upon 10 core axioms (AxiomA--K) formalized in Lean 4: AxiomA (weaving structure), AxiomB (causal partial order), AxiomC (quantum amplitude), AxiomD (operational weaving), AxiomE (information capacity), AxiomF (continuum limit), AxiomG (quantum geometry coupling), AxiomH (gauge group embedding), AxiomI (information causality), AxiomJ (dynamic evolution), and AxiomK (eternal present/scale dynamics). All axioms and proofs are implemented and verified in Lean 4 using the mathlib library. The complete codebase comprises 50 Lean files with approximately 24,300 lines of formal proof, with 2196 compilation tasks passing with zero errors.

\section{Duality Two-One Theorem}

The central structural theorem of CSQIT is the Duality Two-One Theorem, which establishes a discrete complementarity principle deeper than Heisenberg uncertainty:

\begin{theorem}[Duality Two-One]
In the standard framework, if the output function is non-degenerate (causal non-triviality), then the amplitude function is degenerate (information degeneracy). Conversely, if the amplitude function is injective (information non-triviality), then the output function is constant (causal degeneracy).
\end{theorem}

\textit{Proof sketch}. The proof follows from four elementary observations: (1) amplitude is a semigroup homomorphism into $\mathbb{C}$; (2) unitarity ($|\alpha| = 1$) ensures $\text{amplitude}(\alpha) \neq 0$; (3) in a finite set, injectivity of a self-map implies surjectivity; (4) AxiomA's \texttt{compose\_output} then forces output degeneracy. This establishes discrete complementarity---a structural prohibition against simultaneous causal and informational non-triviality.

Fin 7 resolves this tension by relaxing the standard \texttt{compose\_output} constraint, enabling both aspects to coexist non-trivially through a semigroup structure.

\section{Characteristic Constant $\theta$}

$\mathbb{Z}_7$ is the smallest prime-order cyclic group satisfying both unitarity and injectivity of the amplitude function. For this model, the characteristic constant is:

\[
\theta = \frac{1}{2 + 2\cos(2\pi/7)} \approx 0.308
\]

This constant satisfies the cubic equation $\theta^3 - 6\theta^2 + 5\theta - 1 = 0$. Under the duality interpretation, $\theta$ corresponds to the observed total matter density $\Omega_m = 0.311$ (Planck 2018)~\cite{Planck2018} with a deviation of approximately $1\%$. The theoretical value 0.308 lies within the Planck 2018 1$\sigma$ confidence interval $[0.305, 0.317]$.

The cubic nature of $\theta$ is not incidental: the algebraic number field extension $\mathbb{Q}(\zeta_7 + \zeta_7^{-1})$ has degree $(7-1)/2 = 3$, making $\mathbb{Z}_7$ the minimal structure supporting irreducible three-way weaving flow.

\section{Uniqueness of $p=7$}

The algebraic expansion spectrum $\theta(p) = 1/(2+2\cos(2\pi/p))$ is strictly monotonically decreasing for primes $p \ge 3$:

\begin{table}[h]
\centering
\begin{tabular}{ccc}
\hline
$p$ & $\theta(p)$ & In $(0.28,0.33)$? \\
\hline
3 & 1.000 & No \\
5 & 0.382 & No \\
\textbf{7} & \textbf{0.308} & \textbf{Yes} \\
11 & 0.272 & No \\
$\infty$ & 0.250 & No \\
\hline
\end{tabular}
\caption{Algebraic expansion spectrum. Only $p=7$ yields $\theta(p)$ within the structure formation window.}
\label{tab:spectrum}
\end{table}

The window $(0.28, 0.33)$ is not arbitrary: $\Omega_m < 0.28$ prevents gravitational collapse into galaxies; $\Omega_m > 0.33$ causes premature recollapse before structure formation. The strict monotonicity of $\theta(p)$ ensures that among all primes, only $p=7$ occupies this window---a direct constraint from algebraic number theory on which causal lattices can support complex structure.

\section{Existential Inversion}

We reframe the question ``why 7?'' as an epistemological inversion:

\begin{quote}
``Structure$=7$ is the necessary condition for observers who can ask `why 7' to exist.''
\end{quote}

Formally, let $S(p)$ denote the statement ``a causal lattice with algebraic base $p$ can support internal observers.'' Then $S(p) \Rightarrow p = 7$. The justification is structural:
\begin{itemize}
    \item For $p=3$: $\theta=1.0$, absolute closure---no evolution, no time arrow, no observers.
    \item For $p=5$: $\theta=0.382$, quadratic recurrence---reversible oscillation, no irreversible records, no observers who can ``remember'' why.
    \item For $p \ge 11$: $\theta < 0.272$, density too low---no bound structures, no galaxies, no observers.
    \item Only $p=7$ provides the cubic nonlinearity enabling past-present-future closure, irreversible records, and self-referential cognition.
\end{itemize}

This is not the weak anthropic principle (``we happen to be here'') but a strong logical-closure theorem (``only here allows `here' to be defined''). The reason we can ask ``why 7?'' is precisely because we are structural products of 7.

\section{Discussion}

CSQIT demonstrates that a formalized axiomatic approach---with zero free parameters---can produce quantitative results comparable to observational data. The proximity of $\theta \approx 0.308$ to $\Omega_m = 0.311$ is the first observational anchor of this structural deduction.

The complete source code, including all Lean 4 proofs, is available at \url{https://github.com/New-Beginning-Universe-Research-Group/CSQIT}.

\acknowledgments

This work was conducted independently without external funding.

\bibliography{references}

\end{document}
